\chapter*{Conclusion générale et perspectives}
\addcontentsline{toc}{chapter}{Conclusion générale et perspectives}

\section*{Bilan du travail réalisé}

Ce stage avait pour objectif de concevoir, réaliser et évaluer un assistant de recherche
juridique souverain appliqué au Code du travail marocain. Le système livré répond à une question
posée en langage naturel en s'appuyant uniquement sur les articles réellement retrouvés dans le
corpus, cite systématiquement leurs numéros, et s'abstient lorsque l'information n'y figure pas.
L'ensemble de la chaîne s'exécute en local, sans aucun appel à une API externe : la question du
citoyen ne quitte à aucun moment la machine.

Les trois missions engagées ont été menées à terme :

\begin{itemize}
    \item \textbf{Le corpus} compte 589 articles structurés, enrichis de leurs métadonnées de
    hiérarchie, et se régénère à l'identique depuis le PDF source --- propriété vérifiée octet
    pour octet depuis un clone vierge du dépôt.
    \item \textbf{Le moteur RAG} combine recherche sémantique et lexicale, encadrée par deux
    garde-fous indépendants du modèle : un seuil d'abstention appliqué avant l'appel au LLM, et
    une vérification systématique des citations produites.
    \item \textbf{L'évaluation} repose sur un jeu de référence de 37 questions et trois scripts
    rejouables. Aucun chiffre du rapport n'a été saisi à la main.
\end{itemize}

\section*{Principaux résultats}

Sur le jeu de référence, la recherche atteint un Recall@5 de 0{,}97 et un MRR de 0{,}89. Côté
réponses, 97~\% citent au moins un article, 94~\% des citations sont vérifiées comme réellement
fournies au modèle, et l'abstention est correcte sur la totalité des questions hors périmètre.
Le temps de réponse médian est de 3{,}5~secondes, réduit à 0{,}14~seconde lorsque le garde-fou
se déclenche.

Deux résultats méritent d'être soulignés parce qu'ils contredisent les attentes initiales :

\begin{enumerate}
    \item \textbf{La recherche dense seule dépasse la recherche hybride} sur ce jeu de questions.
    Les questions en langage naturel étant des reformulations, elles favorisent le sémantique ;
    la fusion RRF dilue alors un bon classement dense.
    \item \textbf{Une hallucination inter-codes, invisible aux tests manuels}, a été mise au jour
    par l'évaluation : le modèle refusait correctement une question hors périmètre, puis inventait
    une réponse citant une loi inexistante. Sa correction, puis le rééquilibrage de cette
    correction, n'ont été possibles que parce qu'un banc de mesure rejouable existait.
\end{enumerate}

Ces deux constats justifient rétrospectivement le poids donné à l'évaluation dans ce projet : ils
n'auraient été accessibles ni par l'intuition, ni par des tests ponctuels.

\section*{Limites}

Plusieurs limites sont assumées et documentées :

\begin{itemize}
    \item \textbf{La vérification des citations est syntaxique, non sémantique.} Elle établit
    qu'un article cité figurait dans le contexte, pas qu'il dise ce que la réponse lui attribue.
    \item \textbf{Trois questions valides sur 32 sont refusées à tort}, coût assumé d'une
    abstention volontairement stricte.
    \item \textbf{Le périmètre reste un code unique} et une seule langue.
    \item \textbf{Le corpus repose sur une version datée} du texte : sa mise à jour reste un acte
    humain, déclenché et validé, non une veille automatisée.
\end{itemize}

\section*{Perspectives}

\paragraph{Extension à un second code.}
L'architecture a été conçue pour généraliser, mais une évolution structurante reste nécessaire :
les numéros d'articles cessent d'être uniques dès qu'un second code entre dans le corpus ---
« l'article 32 » existe dans plusieurs codes. Chaque article devra donc porter un identifiant
composite associant le code et le numéro. Le Code de la famille constitue le candidat naturel,
étant cité dans le sujet de stage et parmi les plus consultés par le citoyen. Sa sensibilité
particulière imposera une discipline d'abstention encore plus stricte.

\paragraph{Couche multilingue.}
L'approche retenue consisterait à traiter l'arabe et l'anglais comme une couche de traduction
greffée aux extrémités du système, sans toucher au moteur : la question est ramenée au français,
le pipeline s'exécute inchangé, et la réponse est retraduite --- les numéros d'articles cités
restant, eux, indépendants de la langue. Le darija, dépourvu d'orthographe standardisée et
sujet à une forte alternance codique, est délibérément différé.

\paragraph{Reclassement par cross-encodeur.}
Prévu comme option conditionnelle, il reste à mesurer. Compte tenu du constat sur la supériorité
du dense, son apport devrait être évalué contre cette configuration et non contre l'hybride.

\paragraph{Vérification sémantique des citations.}
C'est la limite la plus intéressante à lever : vérifier non plus la présence de l'article cité,
mais la cohérence entre l'affirmation produite et le contenu réel de cet article.

\section*{Apport personnel}

Ce stage m'a permis de mettre en pratique une chaîne complète de traitement automatique du
langage, de l'extraction d'un corpus brut jusqu'à l'évaluation quantitative d'un système
génératif. Au-delà des compétences techniques --- recherche sémantique, bases vectorielles,
modèles de langage locaux --- il m'a surtout appris la valeur d'une démarche de mesure
systématique : plusieurs des décisions défendues dans ce rapport contredisent l'intuition que
j'en avais au départ, et n'ont pu être tranchées que par des chiffres reproductibles.
